\section{Experimental Design for the Extensions}\label{sec:methodology}
\subsection{What Is Tested, and on Which Programs?}
The evaluation separates three practical questions: whether Zone identifies
embedded data correctly, whether a transformed program still performs its
expected task, and how often static scanners flag the resulting file.
Table~\ref{tab:test_populations} summarizes the inputs used for each question.

\begin{table}[t]
\caption{Test populations. The thirteen controlled programs and the thirteen
selected x64 payloads are different sets with different purposes.}
\label{tab:test_populations}
\centering\small
\begin{tabularx}{\linewidth}{@{}l r Y@{}}
\toprule
Population & Inputs & Purpose \\
\midrule
Controlled functional programs & 13 & Check exact expected output after transformation. \\
Annotated classification programs & 6 & Check which bytes are identified as data. \\
Original Windows payload corpus & 102 & Measure construction, functionality, and static detection. \\
Recovery subset of that corpus & 49 & Compare historical and updated working rates on previously problematic inputs. \\
Data-handling subset of that corpus & 19 & Compare symbolic and LLM identification within the same pipeline. \\
\bottomrule
\end{tabularx}
\end{table}

The thirteen \emph{controlled functional programs}, also called test
fixtures, are small purpose-built assembly programs: six cases on each
architecture cover arithmetic, conditional branches, loops, local calls,
combined control flow, and embedded strings. One additional x64 case
exercises access to embedded data through a computed code address.
These programs are 19--68 bytes long. Each writes a known result into a
256-byte output buffer; a separate native C runner executes the program
and compares the whole buffer with the expected bytes. Untouched positions
must retain their initial value. All thirteen originals pass this check.
This compact suite tests individual mechanisms; the original payload corpus
provides the broader application-level evaluation.

The six \emph{annotated classification programs} answer a different
question: does the method locate the data correctly? Two contain multiple
strings (one per architecture), one x64 case contains a binary table,
two contain no separate data, and one x64 case contains instruction bytes
that resemble text. For each program the correct data locations are written
down before either classifier runs. The model does not receive these answers.
The two suites are complementary: one checks the result of execution,
the other checks the identification decision itself.

The original corpus contains 38 x64 and 64 x86 Windows payloads.
Functional tests check the expected dialog, network connection event,
or clean termination for the relevant payload class. All 102 originals pass.
Throughout this article, \emph{working} means that the transformed output
passes the corresponding functional test. We use this empirical definition
consistently for the historical and extended evaluations.

The recovery comparison selects the 49 original inputs whose historical
working rate was below 50\%: thirteen x64 and 36 x86. It compares the same
inputs and fifteen transformation configurations before and after the
pipeline update. New results use ten variants per x64 input/configuration
and five per x86 input/configuration. Data-preparation-only controls are
reported separately. The comparison measures the updated pipeline as a
whole, including the improved MessageBox test; it is not an ablation of
data identification alone.

\subsection{Configurations, Repetitions, and Model Calls}
An experimental \emph{cell} is one combination of input, transformation,
and prompting condition. For example, one x64 test program transformed
with EQ using a prompt with two examples is one cell. Running that
combination three times gives three repetitions of the same cell.

The native controlled benchmark uses three runs per input and transformation
family. The LLM benchmark uses four families, two prompting conditions,
and three separate single-generation requests per cell. This gives
$7\times4\times2\times3=168$ x64 requests and
$6\times4\times2\times3=144$ x86 requests. All rejected responses and
transport failures remain in the counts. The classification benchmark uses
three LLM requests and one symbolic analysis per annotated program.

On the full original corpus, wrapping uses three randomization seeds per
input, giving $102\times3=306$ planned variants. The LLM experiment has
$102\times4\times2=816$ cells and an output-token limit of 65,536.
It targets one accepted output per cell; two successful cells required a
second generation, and fourteen failed cells were restarted without success.
These results describe the completed campaign, including those retries.

The data-handling comparison uses nineteen original inputs (thirteen x64,
six x86) with data extracted by symbolic preparation. Each method proposes
data regions, after which the same deterministic handling and functional
test are applied. Input selection and shared processing are thus held
explicit in interpreting the comparison.

The accompanying artifact records engine versions and request histories.
Native x86 IBP was tested with an earlier implementation than the later
x86 LLM and wrapping experiments; these are separately identified in the
artifact and are not a controlled comparison of engine versions.

\subsection{Build Rate, Working Rate, and Classification Accuracy}
Let $N$ be planned cases, $B$ produced binaries, and $W$ binaries passing
the functional test. Build rate is $B/N$, overall working rate is $W/N$,
and working rate among produced binaries is $W/B$.
For example, ten planned cases, eight produced binaries, and six successful
tests give 80\%, 60\%, and 75\%, respectively. The recovery figure compares
successful tests with recorded tests, matching the available historical
records. A repeated execution of the same binary is not a new variant.

For classification, \emph{predicted} means the bytes that the method labels
as data. \emph{Reference} means the known correct data bytes annotated
before the test. For example, if a string occupies bytes 20--29, the
reference contains those ten byte positions. A prediction of bytes 20--27
finds eight of them and introduces no extra byte: precision is 100\%
and recall is 80\%. In general, for predicted positions $P$ and reference
positions $G$, precision is $|P\cap G|/|P|$ and recall is $|P\cap G|/|G|$.
An exact match requires all data intervals to agree, with no omissions or
additions. For programs without separate data, the correct answer is an
empty set. On original payloads without independent annotations, we report
agreement between the methods and the resulting execution outcomes.

\subsection{Similarity, Detection Rate, Diversity, and Cost}
\paragraph{Similarity.}
The LLM similarity score compares the prepared input assembly with the
accepted output assembly. It measures how much of their instruction
sequences matches, using $S=2M/(u+v)$ for $M$ matched elements and sequence
lengths $u,v$~\cite{python_difflib}. NOP-like operations are ignored for
the NOP condition, so adding only those operations can leave $S=1$ even
when the binary changes. Similarity describes representation; functionality
is measured by the execution tests. This score is separate from the
historical cosine and entropy measures.

\paragraph{Detection rate.}
We retain the historical engine-level detection rate:
$d_i=m_i/a_i$, where $m_i$ engines flag file $i$ as malicious and $a_i$
engine responses are reported. The tables give its mean across files,
$\overline d$, as a percentage. We also report the \emph{detected-file
rate}, $q$: the fraction of files flagged by at least one engine.
For example, one malicious response out of 75 is an engine detection rate
of 1.33\%; that file still counts as detected in $q$. Thus the two rates
answer different questions: how many engines flag a file, and how many
files receive any malicious verdict.

Nonempty current reports contain 75 responses, including unsupported
formats, timeouts, and failures. We disclose these categories separately;
an empty report is treated as missing. These measurements concern static
scans of raw files. Passing and failing outputs remain separately identified.

\paragraph{Diversity and resources.}
$U$ counts distinct output SHA-256 hashes: identical output bytes count
once for diversity and scanning, even when produced by multiple runs.
A shared hash retains membership in every relevant experimental condition.
Consequently, per-condition unique counts need not sum to the global count.
Time is the recorded duration of a run, including local processing.
Estimated API costs include failed calls and restarted runs. Source records,
version mappings, and the detailed accounting are supplied with the artifact.
